\section{Hamiltonian and Optimization Details}
\label{app:optimization}

\subsection{Statewise Hamiltonian evaluation}

For edge-bit row vector $x$, edge-cost vector $c$, resource vector $r$, edge--
vertex incidence matrix $A$, and demand vector $d$, the implementation computes
\begin{equation}
 C(x)=xc,qquad
 f_{\mathrm{raw}}(x)=\|xA-d\|_2^2,qquad
 u(x)=\max\{0,xr-B\}.
\end{equation}
The scale-controlled functions and normalized energy then follow the
controlled-energy definition in the main article. Arrays are evaluated in deterministic chunks;
the diagonal energy vector has length $2^m$ and is not stored outside the
active simulation run.

The controlled penalty contract was audited before QAOA. Both the historical
raw and prospective controlled contracts returned an exact original RCSP
optimum on 140/140 tasks. The controlled absolute span was 490.51--744.17,
compared with $3.42698\times10^7$--$1.82405\times10^9$ for the historical raw
form. Normalization by 172 gives pilot spans 2.97815--4.32657. Ground-state
correctness does not imply perfect phase-scale invariance; the remaining
within-base relative range has median 6.8\% in the selected pilot tasks.

\subsection{Exact statevector circuit}

The initial vector has amplitude $2^{-m/2}$ on every basis state. A cost layer
multiplies amplitude $x$ by $e^{-i\gamma H_C(x)}$. Each factor
$e^{-i\beta X_e}$ is applied by pairwise amplitude updates on bit $e$, so a
full mixer implements $e^{-i\beta\sum_eX_e}$. Parameter arrays store all
$\gamma$ angles followed by all $\beta$ angles. State norm is checked after
simulation; no raw statevector is persisted.

All arms use the same $H_C$ for the cost phase. O1, O2, and O3 do not replace
the quantum Hamiltonian; they replace only the scalar classical loss computed
from the terminal probability distribution. Exact weighted CVaR sorts the
energy/probability pairs stably, accumulates probability until $\alpha=0.10$,
and uses only the fractional mass necessary at the cutoff.

\subsection{COBYLA accounting}

COBYLA is called with initial radius 0.5 and constraint tolerance $10^{-8}$.
Pilot random starts sample $\gamma$ in $[0,2\pi]$ and $\beta$ in
$[0,\pi]$. These are initialization domains only: subsequent optimizer
coordinates are unbounded, and the implementation applies neither a bound nor
modulo-$2\pi$ wrapping to $\gamma$. Because the normalized Hamiltonian
generally has a noninteger spectrum, no fundamental $2\pi$ periodicity in
$\gamma$ is assumed. The pilot uses 120 objective evaluations per
independently initialized run, continuation uses 120 from the embedded point,
and the matched objective and held-out $p=3$ arms use 240. Diagnostic start and
terminal evaluations remain separate from the optimizer's counted calls.

For the pilot's primary metric view, each metric is summarized by its median
over three fixed seeds. A secondary multistart view selects the seed only by
minimum terminal mean energy and subsequently reports that seed's operational
metrics. Feasibility and optimality never participate in seed selection. The
objective-comparison arms share the selected $p=2$ initialization.

\section{Optimizer Attribution}
\label{app:optimizer-attribution}

\subsection{Nested identity and certified failures}

The exact $p=2\to3$ embedding defined in the main article was evaluated for every one of the
56 tasks and three $p=2$ seeds. Maximum statevector, probability, and objective
differences were all zero at stored precision. The original $p=3$ terminal
energy exceeded the embedded point by more than $10^{-10}$ in 29 of 168 runs.
The median signed gap over all runs was $-0.199582$; this overall median does
not diminish the certification of the 29 positive-gap cases.

Continuation improved the mean objective for 29/29 certified failures and
$\gfeas$ for 27/29. Across all 168 runs, continuation classifications were
151 ``objective and feasibility improve'' and 17 ``objective improves,
feasibility worsens.'' Terminal COBYLA points matched the tracked best
evaluated objective within $10^{-10}$ in every continuation run.

\subsection{Budget and optimizer controls}

For one objective-selected seed per task, independent continuations at 120,
240, and 480 evaluations tested budget response. Median objective gains were
0.03197 from 120 to 240 and 0.01386 from 240 to 480; corresponding median
$\gfeas$ changes were 0.08196 and 0.02410. Objective improved while
$\gfeas$ worsened in 4/56 and 10/56 transitions, respectively. Thus additional
budget did not remove objective--feasibility non-equivalence.

A matched 120-evaluation Nelder--Mead control had a median objective difference
of $-8.95\times10^{-5}$ relative to COBYLA and a median $\gfeas$ difference of
$-0.1168$. It reached lower energy on 29/56 tasks; 15 of those also had lower
$\gfeas$. This control is diagnostic, not an optimizer leaderboard, because
only one configuration per method was tested.

These early controls are retained as part of the frozen attribution record.
The later reviewer-robustness matrix in \cref{app:reviewer-robustness} uses
all three original seeds, matched actual objective-call caps, three local
optimizers, and a separately frozen $p\leq4$ depth--budget subset; it
supplements rather than replaces these historical results.

\subsection{Decomposition and response slices}

For every matched state, the stored mean objective equals the sum of routing,
flow-penalty, and resource-penalty components within numerical tolerance. The
four mutually exclusive mass categories---valid and resource-feasible,
flow-invalid/resource-feasible, flow-valid/resource-violating, and both
violating---sum to one. The medians reported in the main article's
optimizer-attribution section are
computed from seed-level paired differences, not differences of separately
aggregated medians.

Fifteen deterministic two-dimensional response slices varied the added
$(\gamma_3,\beta_3)$ while freezing earlier layers. On a $41\times41$ grid,
median objective gain over the embedded point was 0.05648 and median maximum
$\gfeas$ gain was 0.33887. These slices visualize available local response but
cannot certify the global six-parameter optimum.
